\chapter{Additional Validation}
\label{app:additional-validation}

\section{Scaled chi-squared audit design}

The component degree of freedom in Equation~\eqref{eq:nu-p-hat} is based on a
scaled chi-squared working model for the estimated MI variance. A separate
simulation isolated this mechanism from the final two-sample test. It used 64
fixed scenarios and evaluated both population components, giving 128
component distributions. For each component, 10,000 multinomial tables
estimated oracle moments and an independent 10,000 tables evaluated fit.

The audit made two comparisons. The oracle comparison supplied the empirical
mean and variance to scaled chi-squared, normal, and lognormal distributions.
Differences between these models therefore measure distributional shape. The
population first-order comparison instead used $V(P)$ and
$\tau^2(P)/n_P$ from the derivation, so it includes error in the predicted
moments.

Kolmogorov--Smirnov (KS) distance measures whole-distribution disagreement.
Tail error is the absolute difference between the observed exceedance
probability at a candidate quantile and its nominal probability.

\begin{table}[htbp]
\centering
\caption{Overall distributional audit for the MI variance estimator.}
\label{tab:appendix-scaled-chi-overall}
\small
\resizebox{\textwidth}{!}{%
\begin{tabular}{@{}lccccc@{}}
\toprule
Model & Mean KS & Median KS & Tail error .05 & Tail error .01 & Win rate \\
\midrule
Oracle scaled chi-squared & 0.0199 & 0.0143 & 0.0045 & 0.0026 & 0.4219 \\
Oracle normal & 0.0196 & 0.0143 & 0.0054 & 0.0038 & 0.4688 \\
Oracle lognormal & 0.0291 & 0.0227 & 0.0062 & 0.0040 & 0.1094 \\
Population first-order chi-squared & 0.1622 & 0.0822 & 0.0765 & 0.0345 & -- \\
\bottomrule
\end{tabular}%
}
\end{table}

The oracle normal model had the smallest mean KS distance by 0.0003. The
scaled chi-squared model had smaller upper-tail errors at both evaluated
levels. The chi-squared working family is therefore competitive for the part
of the distribution that determines two-sided rejection, although it is not
uniformly the closest family.

\section{Audit by table shape}

\begin{table}[htbp]
\centering
\caption{Oracle distributional fit by table shape.}
\label{tab:appendix-scaled-chi-shape}
\small
\resizebox{\textwidth}{!}{%
\begin{tabular}{@{}>{$}c<{$}cccccc@{}}
\toprule
\text{Shape} & Chi KS & Normal KS & Chi tail .05 & Normal tail .05 & Empirical skew & DF ratio \\
\midrule
(2\times2) & 0.0354 & 0.0273 & 0.0080 & 0.0075 & 0.1920 & 0.9895 \\
(3\times3) & 0.0209 & 0.0197 & 0.0056 & 0.0058 & 0.2095 & 0.9695 \\
(5\times5) & 0.0128 & 0.0176 & 0.0026 & 0.0044 & 0.1827 & 0.9658 \\
(10\times10) & 0.0105 & 0.0137 & 0.0017 & 0.0038 & 0.1071 & 0.9568 \\
\bottomrule
\end{tabular}%
}
\end{table}

The scaled chi-squared shape becomes more competitive as table dimension
increases. The final column is the ratio of the plug-in component degree of
freedom to the empirical moment-matched degree of freedom. The median ratio
over all 128 components was 0.9703.

\section{Audit by sampling regime}

\begin{table}[htbp]
\centering
\caption{First-order and oracle scaled chi-squared diagnostics by regime.}
\label{tab:appendix-scaled-chi-regime}
\small
\resizebox{\textwidth}{!}{%
\begin{tabular}{@{}lrrrrr@{}}
\toprule
Regime & Components & DF ratio & Oracle KS & First-order KS & Tail .05 \\
\midrule
Well sampled & 16 & 0.9585 & 0.0166 & 0.2172 & 0.0960 \\
Moderate & 16 & 0.9546 & 0.0231 & 0.2065 & 0.0963 \\
Sparse and imbalanced & 16 & 0.9649 & 0.0188 & 0.1488 & 0.0622 \\
Highly skewed and sparse & 16 & 0.9946 & 0.0164 & 0.0462 & 0.0128 \\
Ultra-skewed and sparse & 16 & 1.0036 & 0.0173 & 0.0786 & 0.0150 \\
Widespread sparsity & 16 & 0.8522 & 0.0272 & 0.3081 & 0.2042 \\
Equal-MI shape mismatch & 16 & 0.9658 & 0.0236 & 0.1437 & 0.0609 \\
Extreme sample imbalance & 16 & 0.9881 & 0.0161 & 0.1483 & 0.0648 \\
\bottomrule
\end{tabular}%
}
\end{table}

The largest failure occurs under widespread sparsity. The oracle
chi-squared fit also deteriorates there, and first-order moment error becomes
substantial. This agrees with the lower valid rate and residual calibration
error in the main experiment.

\section{Interpretive limit}

This component audit does not make the final Student reference exact. The MI
contrast and its estimated standard error are computed from the same tables
and can remain dependent. The audit establishes that scaled chi-squared
moment matching is a plausible model for the positive variance estimator,
especially in the upper tail, while identifying the first-order moments as a
larger source of finite-sample error.
